\section*{Appendix B: How to Use This Method in Practice}\label{section-practical-guide}

This section provides step-by-step guidance for applied researchers on implementing the noise decomposition framework.

\subsection*{Step 1: Prepare the Response Matrix}

Organize your data as a three-dimensional array (or long-format data frame):
\begin{itemize}
    \item \textbf{Cases (i):} Rows represent distinct cases, patients, or items to be rated
    \item \textbf{Raters (j):} Columns represent individual raters or experts
    \item \textbf{Occasions (t):} If available, include repeated ratings of the same case by the same rater
    \item \textbf{Missing data:} The Bayesian framework handles missing data naturally; not all raters need to rate all cases
\end{itemize}

\subsection*{Step 2: Choose Outcome Coding}

For binary outcomes:
\begin{itemize}
    \item Code positive/yes/present as 1, negative/no/absent as 0
    \item \textbf{Handling "unsure" or "indeterminate":} Common approaches include:
    \begin{enumerate}
        \item Collapse "unsure" to 0 (conservative for diagnosis)
        \item Collapse "unsure" to 1 (sensitive for screening)
        \item Treat as missing data (if infrequent)
        \item Use an ordinal model (see Extensions section)
    \end{enumerate}
    Document your choice and justify it based on the clinical context.
\end{itemize}

\subsection*{Step 3: Fit the Model}

\subsubsection*{Using PyMC (Python):}
\begin{verbatim}
import pymc as pm
import numpy as np

# Assume data: Y[i,j,t] is response matrix
# N cases, M raters, T occasions

with pm.Model() as model:
    # Priors on standard deviations
    sigma_c = pm.HalfNormal('sigma_c', sigma=2)
    sigma_l = pm.HalfNormal('sigma_l', sigma=2)
    sigma_p = pm.HalfNormal('sigma_p', sigma=2)
    sigma_o = pm.HalfNormal('sigma_o', sigma=1)

    # Non-centered parameterization
    z_c = pm.Normal('z_c', 0, 1, shape=N)
    z_l = pm.Normal('z_l', 0, 1, shape=M)
    z_p = pm.Normal('z_p', 0, 1, shape=(N,M))
    z_o = pm.Normal('z_o', 0, 1, shape=(N,M,T))

    c = pm.Deterministic('c', sigma_c * z_c)
    l = pm.Deterministic('l', sigma_l * z_l)
    p = pm.Deterministic('p', sigma_p * z_p)
    o = pm.Deterministic('o', sigma_o * z_o)

    # Intercept
    ybar = pm.Normal('ybar', 0, 2)

    # Linear predictor
    z = c[i_idx] + l[j_idx] + p[i_idx,j_idx] + o[i_idx,j_idx,t_idx] + ybar

    # Likelihood
    Y_obs = pm.Bernoulli('Y_obs', logit_p=z, observed=Y_flat)

    # Sample
    trace = pm.sample(2000, tune=2000, chains=4, target_accept=0.98)
\end{verbatim}

\subsubsection*{Using R (lme4 for approximate ML):}
For a simplified model without occasion noise:
\begin{verbatim}
library(lme4)
# Data in long format with columns: case_id, rater_id, response
fit <- glmer(response ~ (1|case_id) + (1|rater_id) + (1|case_id:rater_id),
             data=dat, family=binomial(link="logit"))
\end{verbatim}

Note: \texttt{lme4} provides point estimates; for full uncertainty quantification, use Bayesian methods.

\subsection*{Step 4: Check Convergence and Model Fit}

\textbf{Convergence diagnostics (Bayesian):}
\begin{itemize}
    \item $\hat{R} < 1.01$ for all parameters
    \item Effective sample size $> 400$ per chain
    \item Visual inspection of trace plots (should resemble "fuzzy caterpillars")
\end{itemize}

\textbf{Posterior predictive checks:}
\begin{itemize}
    \item Compare observed vs. predicted response frequencies by case
    \item Compare observed vs. predicted response frequencies by rater
    \item If $T \geq 2$: Compare observed vs. predicted switch rates between occasions
    \item Examine residual patterns for systematic deviations
\end{itemize}

\subsection*{Step 5: Report Results}

\textbf{Minimal reporting checklist:}
\begin{table}[h!]
\centering
\small
\begin{tabular}{p{0.4\textwidth}p{0.5\textwidth}}
\toprule
\textbf{Item} & \textbf{What to Report} \\
\midrule
Sample sizes & $N$ cases, $M$ raters, $T$ occasions, \% missing data \\
Model specification & Link function (logit/probit), which variance components included \\
Priors (if Bayesian) & Distributions and hyperparameters for $\sigma$ parameters \\
Convergence & $\hat{R}$ values, effective sample sizes, convergence plots \\
Variance estimates & Posterior means (or MLEs) and 95\% credible (or confidence) intervals for $\sigma_c, \sigma_\ell, \sigma_p, \sigma_o$ \\
Percent contributions & \% of total variance from each component with uncertainty \\
Model validation & Posterior predictive check results or goodness-of-fit statistics \\
Code availability & Link to GitHub repository or supplementary code \\
\bottomrule
\end{tabular}
\caption{Minimal reporting checklist for noise decomposition studies.}
\label{tab:reporting}
\end{table}

\subsection*{Interpretation Pitfalls to Avoid}

\begin{enumerate}
    \item \textbf{Don't confuse latent-scale and probability-scale variance.} Variance components are on the logit scale. To interpret on the probability scale, use posterior predictive simulation.

    \item \textbf{Don't ignore prevalence effects.} If cases are highly imbalanced (e.g., 90\% positive), variance estimates may be influenced by ceiling/floor effects.

    \item \textbf{Don't assume high agreement means low noise.} A group of raters with similar biases (high level noise, low pattern noise) may show high agreement but poor accuracy.

    \item \textbf{Don't over-interpret small differences in percent contributions.} Check whether credible/confidence intervals overlap before concluding that one noise source dominates.

    \item \textbf{Beware of rater selection bias.} If raters were selected non-randomly (e.g., "problem" raters excluded), estimates of $\sigma_\ell$ and $\sigma_p$ may underestimate true population variability.
\end{enumerate}

\subsection*{Study Design Guidance}

\textbf{How many cases and raters do I need?}

Based on simulation studies: 
\begin{itemize}
    \item \textbf{Minimum for stable estimates:} $N \geq 30$ cases, $M \geq 10$ raters
    \item \textbf{For precise pattern noise estimation:} $N \geq 50$ cases, $M \geq 20$ raters
    \item \textbf{For occasion noise:} $T \geq 2$ occasions for at least 30\% of case-rater pairs
    \item \textbf{Power for group comparisons:} To detect a 20\% difference in $\sigma$ between two rater groups, aim for $M \geq 15$ raters per group
\end{itemize}

\textbf{When is replication (T > 1) worth the cost?}
\begin{itemize}
    \item If occasion noise is a key research question: yes
    \item If raters are scarce but cases are abundant: replicate a subset (e.g., 20\% of cases rated twice)
    \item If costs are low (e.g., automated re-reading): replicate everything
\end{itemize}

